\section{Reproducibility and Artifact Availability}
\label{sec:artifact}

The release candidate contains three public layers: the report sources, bibliography, figures, and claim--evidence matrix; a repository candidate containing the specification, reference runtime, examples, and public tests; and selected conformance summaries with checksums identifying the underlying forensic record.

Evidence Bundle \evidenceid{} has SHA-256 \evidencehash. The repository candidate includes the official MCP summary, negative-control report, A2A report, operating-system matrix, and machine-readable validation summary. Public smoke tests exercise successful execution and receipt verification, secret-egress rejection, duplicate-key rejection, canonical floating-point rejection, and the pure local default policy.

\ifdefempty{\RepositoryURL}{%
Release coordinates are omitted from this private review build and are injected only after separate GitHub and Zenodo publication authorization.%
}{%
The public repository is \url{\RepositoryURL}. \ifdefempty{\ReportDOI}{}{The archived release is identified by DOI \href{https://doi.org/\ReportDOI}{\ReportDOI}.}%
}

\subsection{Reproduction commands}
From the repository root:
\begin{lstlisting}
make test
make technical-report
make audit
\end{lstlisting}
The prepared CI workflow executes the tests and report build on a clean Linux runner. The claim--evidence matrix labels each assertion as executed, implementation-supported, proposed, inherited, red, or blocked.
